% !TEX root = ../main.tex
\section{证明与延迟模型}
\label{app:proofs}

\subsection{渲染局部性}
\label{app:proofs:locality}

提交时，\method 接纳后继轮次的历史，其中第 $t$ 轮的条目是推测条目 $\hat e_t=\langle\varnothing,\hat a_t,\hat o_t\rangle$，而非父轮次的实际条目 $e_t=\langle\vis{t},\act{t},\obs{t}\rangle$。
提交契约保证这两个条目渲染得到的第 $t{+}1$ 轮输入相同。
为了进一步推出它们在所有后续输入中也具有相同的渲染结果，我们采用以下条件。

\begin{enumerate}[label=(R),leftmargin=*,nosep]
\item \textbf{条目局部渲染。}历史渲染得到的文本，是各条目片段的拼接；每个片段由该条目自身内容以及历史中的角色序列确定。
\end{enumerate}

如果聊天模板依据每条消息的内容与角色渲染，并在消息边界放置特殊词元，则满足 (R)；两种模型的模板在移除推理内容后均如此。
相同词元序列解码得到相同文本，因此 $\hat X_{t+1}=\ctx{t+1}$ 蕴含 $\hat e_t$ 与 $e_t$ 生成相同片段。
由于两份历史的角色序列相同，(R) 进一步保证任意后续扩展都会得到相同输入。
如果实现选择在提交时重新追加父轮次的实际条目，就不需要 (R)。

\subsection{\texorpdfstring{\Cref{prop:lossless}}{无损性命题}的证明}
\label{app:proofs:lossless}

设参考执行提交条目 $e_1,e_2,\dots$，并经过状态 $\env{1},\env{2},\dots$。
如果两份历史的角色序列相同，且生成的片段相同，则称其\emph{等价}；由 (R)，等价历史在任意后续扩展下都会渲染为相同输入。
我们对窗口进行归纳，证明每个窗口开始时，\method 的已提交历史都与参考历史等价，且其已提交状态与参考状态相同。
初始情形成立，因为两种执行从相同消息和状态开始。

考虑从第 $t$ 轮开始的窗口。

\paragraph{父轮次。}
根据归纳假设，输入 $\ctx{t}$ 与参考输入相同。
由 (A1)，即使探测、后继轮次和注入草稿并发运行，输出 $Y_t$ 仍与参考输出相同。
若 $\act{t}=\bot$，则提交相同回复，并取消任何后继轮次；由 (A3)，后继轮次没有留下已提交影响，窗口结束时不变式仍成立。
否则，底层机制通过以下两种方式之一提交 $\act{t}$。
如果已派发候选满足 $\hat a_t\equiv\act{t}$，则接纳执行该候选的副本。
这一副本是在第 $t$ 轮期间从已提交状态复制的；由于本轮期间没有其他提交 (A3)，该状态等于 $\env{t}$。
环境对动作的依赖仅通过其规范化形式，因此由 (A2)，得到参考执行对应的 $(\obs{t},\env{t+1})$。
否则，直接在 $\env{t}$ 上执行 $\act{t}$，这就是参考状态转移。
在两种情况下，父轮次结束后的历史均与参考历史等价，状态均等于 $\env{t+1}$。

\paragraph{后继轮次被提交。}
假设已经启动后继轮次，且最终检查时 $\contract_t=1$。
于是 $\hat X_{t+1}=\ctx{t+1}$ 且 $\hat g_{t+1}=g_{t+1}$；同时，由 (A2)，$\phi(\hat S_{t+1})=\phi(\env{t+1})$ 推出 $\hat S_{t+1}=\env{t+1}$。
后继沙箱 $\hat S'$ 是 $\hat S_{t+1}$ 的副本，且只有后继轮次能够修改它。
由 (A1)，后继轮次产生参考输出 $Y_{t+1}$。
后继轮次内部的底层机制与父轮次一样提交调用，只是此时相对于 $\hat S'$ 执行；因此由 (A2)，其观测和执行后状态为参考执行的 $(\obs{t+1},\env{t+2})$。
接纳时，已提交历史成为 $\hist_t\oplus\hat e_t\oplus e_{t+1}$；根据\cref{app:proofs:locality}，它与参考历史等价；已提交状态则成为 $\env{t+2}$。
由 (A3)，没有其他影响进入已提交状态；由 (A4)，接纳时的状态就是检查过的状态。

\paragraph{后继轮次被拒绝或失败。}
如果没有启动后继轮次，或 $\contract_t=0$，或后继轮次失败，其影响均被限制在 $\hat S'$ 及其自身副本内 (A3)，排空后这些副本被丢弃。
随后，在 $\ctx{t+1}$ 和 $\env{t+1}$ 上运行第 $t{+}1$ 轮，父轮次的论证可原样应用于这一轮。

在所有情况下，窗口提交的条目均与参考执行一致，并在参考状态下结束，归纳完成。\qed

\paragraph{交互任务。}
使用模拟用户时，用户是环境的一部分。
窗口不会跨越用户消息，因此只要模拟器对相同对话的响应可复现，上述论证就适用于每个智能体执行片段。

\subsection{\texorpdfstring{\Cref{prop:latency}}{窗口延迟命题}的证明}
\label{app:proofs:latency}

假设锚点正确、$d_t\le R_t$，父轮次和后继轮次的耗时与底层机制下相同，且没有管理开销。
底层机制下，观测在 $\tau_t=d_t+L_t$ 时刻就绪，父轮次在 $\max(R_t,\tau_t)$ 时刻结束，随后第 $t{+}1$ 轮耗时为 $B_{t+1}$：
\[
T_{\text{sub}}=\max(R_t,\tau_t)+B_{t+1}.
\]
若 $\tau_t<R_t$，\method 在 $\tau_t$ 时刻启动后继轮次；当两个轮次都结束时，窗口结束：
\[
T_{\text{TA}}=\max(R_t,\tau_t+B_{t+1}).
\]
于是 $T_{\text{sub}}-T_{\text{TA}}=R_t+B_{t+1}-\max(R_t,\tau_t+B_{t+1})=\min(B_{t+1},R_t-\tau_t)$。
若 $\tau_t\ge R_t$，则不会启动后继轮次，节省时间为 $0=\min(B_{t+1},\slack_t)$。
这就证明了\cref{eq:latency}。

串行执行耗时为 $R_t+L_t+B_{t+1}$。
底层机制节省 $R_t+L_t-\max(R_t,d_t+L_t)=\min(L_t,R_t-d_t)$。
当 $\tau_t<R_t$ 时，\method 节省 $R_t+L_t+B_{t+1}-\max(R_t,d_t+L_t+B_{t+1})=\min(L_t+B_{t+1},R_t-d_t)$。
当 $\tau_t\ge R_t$ 时，两者节省的时间均为 $R_t-d_t$，上述公式仍成立。

对于加速比上界，令 $M=\max(R_t,\tau_t,B_{t+1})$。
两种情况下均有 $T_{\text{TA}}\ge M$，而 $T_{\text{sub}}\le 2M$，因此 $T_{\text{sub}}/T_{\text{TA}}\le 2$。

\subsection{实际收益}
\label{app:proofs:gain}

实际窗口在三个方面偏离理想模型：后继轮次与父轮次竞争 GPU，管理操作需要时间，且部分后继轮次会被拒绝。
令 $A=\max(R_t,\tau_t)$ 和 $B=B_{t+1}$ 分别为底层机制下父轮次与下一轮的耗时，并用撇号表示在 \method 下测得的耗时。

\paragraph{成功提交的窗口。}
设后继轮次在 $r\le A'$ 时刻启动、耗时为 $B'$；等待前的验证耗时为 $v_1$，接纳和清理耗时为 $v_2$。
窗口耗时为 $T_{\text{hit}}=\max(A'+v_1,r+B')+v_2$。
令重叠时间 $O=\min(B',(A'-r)^+)$，开销 $e=\max(A'+v_1,r+B')-\max(A',r+B')+v_2\in[v_2,v_1+v_2]$，则收益为
\[
G_{\text{hit}}=A+B-T_{\text{hit}}=O-(A'-A)-(B'-B)-e.
\]

\paragraph{被拒绝的窗口。}
设排空耗时为 $C_{\text{d}}$，回退轮次耗时为 $B^{\text{f}}$，其他管理开销为 $V$。
于是 $T_{\text{miss}}=A'+C_{\text{d}}+B^{\text{f}}+V$，且
\[
G_{\text{miss}}=-(A'-A)-(B^{\text{f}}-B)-C_{\text{d}}-V.
\]
若父轮次以回复结束，则 $B=B^{\text{f}}=0$，但仍需承担后继轮次造成的资源争用和排空成本。

\paragraph{期望收益。}
使用提交指示变量 $I$，有
\[
\begin{aligned}
\mathbb{E}[G]={}&\mathbb{E}[I\,O]-\mathbb{E}\big[\Delta A+I\,\Delta B+(1-I)\,\Delta B^{\text{f}}\\
&+I\,e+(1-I)(C_{\text{d}}+V)\big],
\end{aligned}
\]
其中 $\Delta A=A'-A$、$\Delta B=B'-B$，且 $\Delta B^{\text{f}}=B^{\text{f}}-B$。
资源争用通过 $\Delta A$ 和 $\Delta B$ 体现；禁用后继轮次内部的草稿复用也会影响这些项。
第一项必须保留联合期望：提交事件与较长重叠时间存在相关性，因此一般有 $\mathbb{E}[I\,O]\neq\Pr(I=1)\,\mathbb{E}[O]$。
如果每个成功提交的窗口至少收益 $g_{\min}>0$，每个其他窗口至多损失 $c_{\max}\ge 0$，那么提交率 $p>c_{\max}/(g_{\min}+c_{\max})$ 足以保证净收益为正。

\subsection{工具延迟执行策略}
\label{app:proofs:deferred}

在延迟执行策略下，后继轮次的全部工具执行，包括其自身的推测派发，都等待提交决策。
若提交，工具在 $\hat S'=\env{t+1}$ 中运行，并由 (A2) 产生参考状态转移，因此\cref{prop:lossless} 仍然成立。
若拒绝，则没有执行过任何后继工具。
由于工具延迟要在提交后支付，节省时间缩减为近似 $\min(G_{t+1},\slack_t)$，其中 $G_{t+1}$ 为后继轮次的生成时间。

\subsection{随机解码}
\label{app:proofs:stochastic}

采用采样时，将 (A1) 替换为以下要求：输出是输入、配置和请求随机种子的确定性函数；同时，后继轮次使用参考执行第 $t{+}1$ 轮原本会使用的种子。
归纳证明即可保持不变。
必须避免的是依赖结果的选择。
对被拒绝的后继轮次重新采样，或按内容、完成时间从多个后继轮次中进行选择，都会改变已提交输出的分布。
\method 不进行上述操作：被拒绝的后继轮次会被丢弃，第 $t{+}1$ 轮只在参考路径上生成一次。
我们的所有实验均采用贪心解码。
